\section{Reglamento y convivencia}

\subsection{Rol del líder}

\begin{enumerate}[label=\alph*)]
    \item El líder (Jose Luis Sepúlveda) será responsable de coordinar las reuniones, distribuir las tareas y velar por el cumplimiento de los plazos establecidos.
    \item El líder centralizará la comunicación con el docente y será el encargado de administrar el repositorio de GitHub del grupo.
    \item El líder podrá delegar tareas específicas, pero la responsabilidad final del cumplimiento del trabajo es compartida por todos los integrantes.
\end{enumerate}

\subsection{Acuerdos de trabajo}

\begin{enumerate}[label=\arabic*.]
    \item \textbf{Reuniones:} el grupo se reunirá, de manera presencial o virtual, al menos una vez por semana durante el desarrollo de cada trabajo, en el día y horario acordado previamente por todos los integrantes.
    \item \textbf{Comunicación:} el canal oficial de comunicación del grupo será un grupo de WhatsApp, donde se informarán avances, dudas y cambios de horario con un mínimo de 24 horas de anticipación.
    \item \textbf{Cumplimiento de plazos:} cada integrante se compromete a entregar su parte del trabajo dentro del plazo interno acordado por el grupo, el cual siempre será anterior a la fecha límite de entrega del docente, para dejar margen de revisión.
    \item \textbf{Calidad del trabajo:} todo contenido entregado debe ser elaborado por el propio integrante, revisado antes de enviarlo, y ajustado a los requerimientos indicados por el líder y el docente.
    \item \textbf{Participación equitativa:} la carga de trabajo se distribuirá de manera equitativa entre los cuatro integrantes, considerando las capacidades y disponibilidad de cada uno.
    \item \textbf{Presentaciones:} todos los integrantes deben estar presentes y preparados para las instancias de presentación oral o defensa del trabajo, salvo justificación debidamente acreditada.
    \item \textbf{Respeto:} las opiniones de todos los integrantes serán escuchadas y consideradas. Las decisiones se tomarán preferentemente por consenso; en caso de desacuerdo, se resolverá por mayoría simple.
\end{enumerate}

\subsection{Método de trabajo en GitHub}

\begin{enumerate}[label=\arabic*.]
    \item El repositorio del grupo se alojará en GitHub bajo la administración del líder, quien agregará como colaboradores a los demás integrantes.
    \item Se utilizará la rama \texttt{main} únicamente para las versiones finales y aprobadas por todo el grupo. El trabajo en desarrollo se realizará en ramas secundarias (por ejemplo, \texttt{feature/nombre-integrante} o \texttt{feature/tarea}).
    \item Cada integrante debe subir sus avances mediante \textit{commits} con mensajes claros y descriptivos que indiquen qué se modificó (ejemplo: \texttt{"Agrega endpoint de autenticación"}).
    \item Antes de fusionar cambios a la rama principal, se recomienda revisar el trabajo mediante un \textit{pull request}, el cual será revisado por el líder u otro integrante antes de aceptarlo.
    \item El repositorio debe mantener organizada la estructura de carpetas: documentación (informe, reglamento), código fuente del proyecto, y recursos adicionales (imágenes, capturas, videos).
    \item Es responsabilidad de cada integrante mantener actualizado su acceso y realizar copias de seguridad locales de su trabajo antes de subirlo.
    \item El historial de \textit{commits} de GitHub podrá ser utilizado como respaldo objetivo para verificar la participación real de cada integrante en caso de conflicto.
\end{enumerate}

\subsection{Resolución de conflictos}

En caso de desacuerdos o conflictos internos que no puedan resolverse mediante diálogo directo entre los integrantes, se seguirá el siguiente procedimiento:

\begin{enumerate}[label=\arabic*.]
    \item Conversación directa entre las partes involucradas, mediada por el líder.
    \item Reunión grupal para exponer el conflicto y buscar una solución consensuada.
    \item De no llegar a acuerdo, se solicitará la intervención del docente de la asignatura como última instancia.
\end{enumerate}
